\documentclass[conference]{IEEEtran}

\usepackage{amsmath,amssymb}
\usepackage{bm}
\usepackage{graphicx}
\usepackage{cite}
\usepackage{hyperref}
\usepackage{algorithm}
\usepackage{algorithmic}
\usepackage{booktabs}
\usepackage{float}



\title{A Hybrid CNN-Quantum Framework for Automated Malignancy Detection in Lymph Node Histopathology Patches}

\author{
\centering
\begin{tabular}[t]{@{}p{0.22\textwidth}@{\hspace{0.015\textwidth}}p{0.22\textwidth}@{\hspace{0.015\textwidth}}p{0.22\textwidth}@{\hspace{0.015\textwidth}}p{0.22\textwidth}@{}}
\centering
{\normalsize Tarakeshwaran S} \\[-2pt]
{\small\itshape Dept. of Computer Science } \\[-2pt]
{\small\itshape and Engineering} \\[-2pt]
{\small\itshape Velammal Engg. College} \\[-2pt]
{\small Chennai, India} \\[-2pt]
{\scriptsize tarakeshwaran.sampath@gmail.com}
&
\centering
{\normalsize Srivaishnav S} \\[-2pt]
{\small\itshape Dept. of Computer Science } \\[-2pt]
{\small\itshape and Engineering} \\[-2pt]
{\small\itshape Velammal Engg. College} \\[-2pt]
{\small Chennai, India} \\[-2pt]
{\scriptsize vaishnav90100@gmail.com}
&
\centering
{\normalsize Pranesh S} \\[-2pt]
{\small\itshape Dept. of Computer Science } \\[-2pt]
{\small\itshape and Engineering} \\[-2pt]
{\small\itshape Velammal Engg. College} \\[-2pt]
{\small Chennai, India} \\[-2pt]
{\scriptsize praneshjd2004@gmail.com}
&
\centering
{\normalsize Dr. Gunasundari S} \\[-2pt]
{\small\itshape Dept. of Computer Science } \\[-2pt]
{\small\itshape and Engineering} \\[-2pt]
{\small\itshape Velammal Engg. College} \\[-2pt]
{\small Chennai, India} \\[-2pt]
{\scriptsize gunasundari@velammal.edu.in}
\end{tabular}
}

\begin{document}
\maketitle

\begin{abstract}
The efficacy of near-term quantum machine learning (QML) is often limited by two primary hurdles: the high cost of representing complex medical images in quantum states and the significant measurement overhead required for similarity-based inference. This paper introduces a hybrid medical-QML framework in which a convolutional backbone with a compact neck first compresses each 96\,$\times$\,96 lymph-node patch into a low-dimensional feature vector, which is then amplitude-encoded and evaluated by a fixed interference circuit \cite{schuld2017distance,schuld2018book}. We employ a fixed interference-based decision rule that relies on the \emph{sign of the real linear overlap} rather than quadratic fidelity, which reduces the measurement shots needed for inference. The learning phase is executed entirely classically through centroid aggregation of training embeddings, which are then used to define static state-preparation unitaries. We characterize the proposed Interference-Based Quantum Classifier (IQC) rigorously: we prove that, under real amplitude encoding, its decision rule is exactly a linear (signed-centroid) classifier and that the quantum circuit reproduces the classical overlap to machine precision, so its genuine near-term advantage is the measurement-efficient \emph{sign} readout rather than the decision function itself. On PatchCamelyon, evaluated on the held-out test split with five-seed statistical validation, the single-prototype IQC attains $82.7\pm1.4\%$, statistically comparable to a fidelity-based quantum kernel and trailing well-tuned classical linear baselines by $\sim\!1.4$ points; a multi-prototype extension surpasses $k$-NN at $K\!\geq\!11$. We further show that the measurement-efficiency claim holds and is governed by the empirical margin, whereas the noise-resilience of the exact-encoding circuit is bounded by its transpiled depth ($L\!\approx\!1268$), with a depolarizing survival threshold near $p\!\approx\!10^{-4}$. These results position the IQC as a measurement-efficient, NISQ-compatible \emph{quantum-executable linear classifier} and clarify where its advantage does and does not lie.
\end{abstract}

\begin{IEEEkeywords}
Quantum Machine Learning, Quantum Classification, Quantum Interference, Linear Overlap, Hybrid Quantum--Classical Systems, Medical Imaging, Histopathology
\end{IEEEkeywords}

\section{Introduction}

\subsection{Background and Motivation}
Despite new quantum primitives, many NISQ classifiers are still constrained by unstable variational training \cite{benedetti2019parameterized,mcclean2018barren} or inference strategies that demand heavy measurement budgets \cite{schuld2019feature,havlicek2019supervised}. These constraints are particularly acute in histopathology, where the high dimensionality of 96\,$\times$\,96 image patches makes direct quantum feature mapping impractical on current hardware \cite{preskill2018nisq}.

Existing QML frameworks frequently rely on iterative parameterized updates that are difficult to scale, or on similarity metrics that necessitate a large measurement budget. In diagnostic imaging, both numerical stability and execution efficiency are paramount for real-world feasibility.

\subsection{Motivation for Interference-Based Classification}
We address this gap with a hybrid medical-QML architecture. A classical CNN backbone with a compact neck compresses each image into a low-dimensional embedding. The embedding is then amplitude-encoded into a quantum state and processed by a fixed interference circuit. Classification is performed by the sign of the real linear overlap between the test state and a class prototype, requiring a small constant number of shots per inference using one fixed circuit per sample. Learning is entirely classical: class-representative prototype states are constructed by centroid aggregation of embeddings, and the quantum processor is used only for inference.

\subsection{Primary Contributions}
The key contributions of this research include:
\begin{itemize}
    \item A \textbf{hybrid medical-QML pipeline} tailored for lymph node histopathology that overcomes quantum encoding bottlenecks via a modular CNN neck for dimensionality reduction.
    \item A \textbf{decision primitive based on linear interference}, utilizing $\mathrm{sign}(\mathrm{Re}\langle\phi|\psi\rangle)$ to establish a linear hyperplane classifier in Hilbert space, together with an \textbf{equivalence analysis} proving that, for real amplitude encoding, this primitive is exactly a signed-centroid linear classifier and that its multi-prototype variant collapses to a single effective hyperplane at inference.
    \item A \textbf{measurement-efficient inference circuit} that operates with a fixed structure and low shot requirements, circumventing iterative variational training; we derive the shot budget from a Hoeffding bound tied to the empirical decision margin and validate it against binomial sampling of the Hadamard test.
    \item A \textbf{classical-to-quantum learning strategy} where class-representative states are aggregated classically, extended to fixed and adaptive multi-prototype regimes that address the single-mode limitation of a single centroid.
    \item A \textbf{rigorous multi-seed evaluation protocol} using the held-out PCam test split (rather than a validation subset) with five independently retrained CNNs, mean\,$\pm$\,std reporting, paired McNemar significance testing, and a physically wired \texttt{qiskit-aer} noise model, which together provide an honest account of where the interference primitive is and is not competitive.
\end{itemize}

\subsection{Structure of the Paper}
The remainder of this paper is organized as follows. Section~II reviews existing systems and related work. Section~III outlines limitations of existing systems. Section~IV presents the proposed interference-based classifier, including the hybrid pipeline, encoding, and decision mechanism. Section~V details the implementation setup, datasets, and evaluation protocol. Section~VI discusses measurement efficiency, noise robustness, limitations, and scalability. Section~VII concludes and outlines future work.

\section{Related Work / Existing Systems}

Quantum machine learning (QML) has progressed from early linear-algebraic subroutines to hardware-focused models tailored to NISQ constraints. We summarize the main approaches and position the linear-interference framework within the remaining gaps.

\subsection{Variational Quantum Methods}
Variational quantum classifiers (VQCs) represent a prominent class of NISQ-era machine learning models \cite{benedetti2019parameterized,abbas2021power}. Variational quantum classifiers use parameterized circuits whose gate angles are tuned during training to minimize a classical objective. While highly adaptable, these models often encounter training bottlenecks, such as the emergence of barren plateaus where gradients can decay exponentially with qubit count \cite{mcclean2018barren}. Furthermore, the high measurement shot requirement for each optimization iteration can place a significant burden on quantum hardware \cite{benedetti2019parameterized}.

\subsection{Quantum Kernel and Metric-Based Strategies}
Alternatively, quantum kernel methods bypass variational training by embedding inputs into quantum feature spaces and computing pairwise overlaps \cite{schuld2019feature,havlicek2019supervised}. Often quantified using fidelity $F(\rho, \sigma)$, these similarity values are utilized as kernel entries in classical algorithms like Support Vector Machines (SVMs). However, these approaches remain measurement-intensive, typically scaling with $O(N)$ for a dataset of size $N$ \cite{schuld2019feature,havlicek2019supervised}. Additionally, as a quadratic observable, fidelity may overlook phase-sensitive information crucial for specific decision boundaries.
Kernel methods shift the burden from training to inference, and the measurement cost of overlap estimation can dominate in practical settings. This motivates exploring alternatives that preserve decision fidelity without requiring extensive pairwise evaluations \cite{schuld2019feature,havlicek2019supervised}.

\subsection{Quantum Geometry and State Discrimination}
Quantum state discrimination (QSD) frames classification as an identification problem among known quantum states \cite{sergioli2019binary,gonzalez2021classification}. This perspective emphasizes the geometric arrangement of states within Hilbert space over circuit parameterization. Similarly, quantum metric learning seeks to optimize a feature map such that inter-class distances are maximized while intra-class variances are minimized. Our work builds upon these geometric insights by introducing a fixed interference circuit designed to extract a linear overlap signal, offering a more streamlined measuring alternative to full tomography or standard fidelity estimation.
The geometric view emphasizes discrimination between quantum states rather than explicit parameter tuning, which aligns with approaches that treat the measurement signal itself as the decision statistic. This perspective supports our use of a fixed interference measurement as a classifier primitive \cite{sergioli2019binary,gonzalez2021classification}.

\subsection{Interference-Based Quantum Classification}

Works like the distance-based classifier by Schuld et al. \cite{schuld2017distance,schuld2018book} pioneered the use of the Hadamard test for pattern recognition by calculating distances between vectors, and kernel-based variants compute signed fidelity sums in superposition \cite{blank2020kernel,blank2022compact}. Our work differs operationally in two ways: (i) the decision variable is the \textbf{sign of the real overlap}, not a fidelity or distance estimate, and (ii) learning uses a classical centroid construction, with a \textbf{fixed interference circuit} used only at inference. This yields a measurement-efficient primitive that avoids variational training and reduces circuit executions per sample.
Interference-based classifiers already demonstrate how overlap signals can encode decision information. Our approach stays within this family but emphasizes a linear overlap test that is aligned with low-shot inference in NISQ settings \cite{schuld2017distance,schuld2018book}.

Table~\ref{tab:comparison} summarizes the comparison of quantum classification paradigms.

\begin{table*}[t]
\centering
\caption{Comparison of Quantum Classification Paradigms}
\label{tab:comparison}
\begin{tabular}{@{}lllll@{}}
\toprule
Feature & Variational (VQC) & Kernel (QSVM) & Fidelity-Based & \textbf{Linear Interference (IQC)} \\ \midrule
Decision Basis & Probabilities $P(y|x)$ & High-dim Kernels & Quadratic Overlap $F$ & \textbf{Linear Overlap $\mathrm{Re}\langle \phi|\psi \rangle$} \\
Learning Mode & In-circuit Optimization & Classical SVM ($O(N^2)$) & Static/Adaptive & \textbf{Static Classical Aggregation} \\
Quantum Circuit & Parameterized (Trainable) & Fixed (Feature Map) & Fixed (Swap Test) & \textbf{Fixed (Hadamard Test)} \\
Shot Complexity & High ($O(M \cdot N)$) & Moderate ($O(N)$) & Moderate ($O(1/\epsilon^2)$) & \textbf{Low (constant shots)} \\
Phase Sensitivity & No & No & No & \textbf{Retains sign of real interference} \\
Optimization Risk & Barren Plateaus & None (in quantum) & None & \textbf{None} \\ \bottomrule
\end{tabular}
\end{table*}

\section{Limitations of Existing Systems}

\subsection{Bottlenecks in VQC and Kernel Paradigms}
Variational models, while versatile, are often hampered by the iterative need for gradient estimation, which remains expensive on current devices. Similarly, kernel-based classifiers, though avoiding variational training, incur significant measurement costs during the inference phase due to repeated pairwise overlap calculations.
These observations suggest that, beyond expressivity, the practical cost of iterative training and measurement overhead is a key factor in near-term viability. This motivates decision primitives that reduce repeated circuit evaluations at inference time and avoid heavy optimization loops where possible \cite{benedetti2019parameterized,mcclean2018barren}.

\subsection{Sign and Sensitivity Issues in Fidelity Measures}
Fidelity-centric approaches typically rely on the $F(\phi, \psi) = |\langle \phi | \psi \rangle|^2$ metric \cite{schuld2019feature}. While robust, this quadratic similarity measurement discards the relative sign of interference terms, potentially limiting the expressiveness of the resulting decision boundaries. These collective inefficiencies—optimization instability, high measurement counts, and the loss of signature phase information—provide the motivation for our proposed linear-interference decision rule.

\section{Proposed Work}

\subsection{Overview of Hybrid Pipeline}
The proposed Interference-Based Quantum Classifier (IQC) follows a modular pipeline: \emph{Input image} $\rightarrow$ \emph{CNN backbone} $\rightarrow$ \emph{compact neck embedding} $\rightarrow$ \emph{quantum state-preparation} $\rightarrow$ \emph{interference-based decision}. The classical module compresses high-dimensional images into an embedding compatible with a small qubit register, while the quantum module performs a fixed interference measurement to obtain a sign-based decision.

Training is executed classically. We first train the CNN backbone and neck on PCam using standard cross-entropy loss to produce discriminative embeddings. After training, the CNN+neck is frozen and used to generate embeddings for all training samples. These embeddings are aggregated into class prototypes (Section~\ref{sec:centroids}), which are compiled into fixed state-preparation unitaries. Quantum hardware is used only for inference: for each test sample we compute its embedding, amplitude-encode it as $|\psi\rangle$, run the fixed interference circuit against the class prototype, and assign a label based on the sign of the interference score estimated over a small constant number of shots.

Fig.~\ref{fig:flowchart} illustrates the complete hybrid classical--quantum pipeline used in the proposed IQC framework.
\begin{figure}[ht]
    \centering
    \includegraphics[width=1.02\linewidth]{flowchart_fig1.png}
    \caption{Flowchart of the hybrid classical-quantum inference pipeline using the Interference-Based Quantum Classifier (IQC).}
    \label{fig:flowchart}
\end{figure}

\subsection{Classical Feature Extraction Module}
The hybrid pipeline described in Section~IV-A begins with a classical CNN backbone that processes 96\,$\times$\,96 RGB lymph node histopathology patches and extracts hierarchical features correlated with malignancy. To make quantum encoding feasible, we insert a compact ``neck'' module that compresses the feature maps into a low-dimensional embedding. The neck acts as a learned projection (analogous to a task-specific PCA), implemented with global average pooling followed by a fully connected layer. The embedding dimension $d$ is chosen to match the available qubit budget via $n = \lceil \log_2 d \rceil$. Fig.~\ref{fig:architecture} provides a visual representation of the backbone and feature neck.

\begin{figure*}[t]
\centering
\includegraphics[width=1.0\linewidth]{architecture_diagram.png}
\caption{Hybrid pipeline: image $\rightarrow$ CNN backbone $\rightarrow$ neck (dimensionality reduction) $\rightarrow$ quantum encoding $\rightarrow$ interference-based classifier tail. The quantum stage uses $n$ data qubits and one ancilla for sign-based inference (small constant shots).}
\label{fig:architecture}
\end{figure*}

\subsubsection{Module-wise Tensor Shapes}
The end-to-end data flow is:
\begin{itemize}
    \item Input image: $(96, 96, 3)$
    \item CNN feature maps: $(h, w, c)$ (architecture-dependent)
    \item Global average pooling: $(c)$
    \item Neck projection: $(d)$
    \item Quantum feature vector (normalized and zero-padded if needed): $(2^n)$
    \item Quantum register: $n$ data qubits + 1 ancilla
    \item Output: 1 interference score (sign used for label)
\end{itemize}

This structure explicitly defines the image $\rightarrow$ neck $\rightarrow$ quantum classifier tail, ensuring the quantum stage operates on a manageable number of qubits while preserving discriminative features.

\subsection{Quantum Encoding Module}
A critical component of the IQC is the explicit mapping of classical feature embeddings into quantum states. We adopt \textbf{amplitude encoding} as the primary feature map $\Phi: \mathbb{R}^d \rightarrow \mathbb{C}^{2^n}$. A classical embedding $\mathbf{x} \in \mathbb{R}^d$ is first normalized and (if needed) zero-padded to length $2^n$, then mapped to an $n$-qubit state via a state-preparation unitary $U(\mathbf{x})$:

\begin{equation}
|\psi\rangle = \sum_{i=0}^{2^n - 1} x_i |i\rangle,
\end{equation}

with $n = \lceil \log_2 d \rceil$. Amplitude encoding provides the highest information density, encoding $d$ features into $n$ qubits. In our implementation, amplitude encoding is realized via Qiskit's \texttt{StatePreparation} gate, which constructs a sequence of uniformly controlled rotations to prepare any normalized complex vector $\mathbf{x}$ as a quantum state $|\psi\rangle$. The interference mechanism itself is independent of the specific compilation used.

We consider three distinct encoding techniques tailored for the IQC architecture:
\begin{enumerate}
    \item \textbf{Amplitude Encoding}: Assigns the normalized components of a classical vector $\mathbf{x} \in \mathbb{R}^d$ to the amplitudes of a $2^n$-dimensional quantum state. While requiring sophisticated state-preparation circuits, this provides optimal information density and ensures that the linear overlap in Hilbert space reflects the Euclidean dot product in the input space.
    \item \textbf{Angle-Based Encoding}: Represents features $x_j$ as rotation angles. While hardware-efficient and shallow in depth, this method introduces trigonometric nonlinearities into the inference signal.
    \item \textbf{Basis-Level Encoding}: Maps classical binary strings to the corresponding computational-basis states of the register. This is computationally simpler but fails to leverage the superposition-driven interference advantage of the IQC.
\end{enumerate}
Our experiments primarily utilize amplitude encoding.

\subsubsection{Linear Separability in High-Dimensional Feature Spaces}
According to Cover's Theorem, a complex pattern classification problem cast in a high-dimensional space nonlinearly is more likely to be linearly separable than in a low-dimensional space \cite{cover1965geometrical}. By mapping data into the $2^n$-dimensional Hilbert space, the IQC leverages this high dimensionality to separate classes using a simple hyperplane.

However, unlike kernel methods that use an implicit feature map, the IQC works with an \textbf{explicit feature map} $\Phi(\mathbf{x})$. The linear overlap $\mathrm{Re}\langle \phi | \psi \rangle$ is computed directly in the state space. The success of the static centroid construction (Section~\ref{sec:centroids}) depends on whether the chosen $\Phi$ places the class distributions in convex regions that can be separated by a single hyperplane.

\subsection{Linear Interference Decision Mechanism}
The central theoretical contribution of this work is the application of linear interference as the main signal for quantum classification \cite{schuld2017distance,schuld2018book}.

\subsubsection{Linear vs. Quadratic Similarity}
In many existing quantum classification frameworks, state-pair similarity is typically assessed through the fidelity metric $F(\phi, \psi) = |\langle \phi | \psi \rangle|^2$ \cite{nielsen2000qcqi}. Despite being a well-grounded measure in quantum information theory, its quadratic character introduces three practical difficulties for NISQ inference:
\begin{enumerate}
    \item \textbf{Sign Loss}: The modulus-square operation $|\cdot|^2$ destroys the sign of the underlying real interference term, erasing discriminative phase information.
    \item \textbf{Non-linear Decision Boundaries}: Boundaries induced by fidelity are second-order surfaces in amplitude space, complicating geometric interpretation.
    \item \textbf{Shot Noise Sensitivity}: Achieving high-precision fidelity estimates demands $O(1/\epsilon^2)$ measurement shots, placing a heavy burden on NISQ resources.
\end{enumerate}

In contrast, the linear overlap $\langle \phi | \psi \rangle$ is complex-valued. Here $|\psi\rangle$ and $|\phi\rangle$ denote quantum states in Hilbert space, and $a_i, b_i$ denote their complex amplitude coefficients. We focus on its real part:
\begin{equation}
s(\psi) = \mathrm{Re}\langle \phi | \psi \rangle.
\end{equation}
where $\mathrm{Re}\langle \phi | \psi \rangle$ denotes the real part of the inner product between the prototype state $|\phi\rangle$ and the input state $|\psi\rangle$.

\subsubsection{Classification Rule and Thresholding}
The classification process hinges on the sign of the real overlap between a test state $|\psi\rangle$ and the class prototype $|\phi\rangle$. The label prediction $\hat{y}$ is formally defined as:
\begin{equation}
\hat{y} = \text{sign}\left( \mathrm{Re}\langle \phi | \psi \rangle \right).
\end{equation}
While a zero threshold implies a balanced boundary, a bias term $b \in \mathbb{R}$ can be integrated to account for asymmetric class priors:
\begin{equation}
\hat{y} = \begin{cases} +1 & \mathrm{Re}\langle \phi | \psi \rangle + b \geq 0 \\ -1 & \mathrm{Re}\langle \phi | \psi \rangle + b < 0 \end{cases}
\end{equation}

\subsubsection{Decision Geometry: The Hilbert Space Hyperplane}
Geometrically, the decision rule $\mathrm{Re}\langle \phi | \psi \rangle = 0$ defines a \textbf{linear hyperplane} that partitions the $2^n$-dimensional Hilbert space into two half-spaces. Fig.~\ref{fig:geometry_comparison} contrasts this linear-interference boundary with the curved, quadratic boundary induced by fidelity-based similarity, and Fig.~\ref{fig:fidelity_geometry} shows the fidelity geometry.

To see this, let $|\psi\rangle$ and $|\phi\rangle$ be represented by their amplitude vectors $\mathbf{c}_\psi, \mathbf{c}_\phi \in \mathbb{C}^{2^n}$. The real part of the inner product can be written as:
\begin{equation}
\mathrm{Re}\langle \phi | \psi \rangle = \frac{1}{2}\left( \langle \phi | \psi \rangle + \langle \psi | \phi \rangle \right) = \text{Re}(\mathbf{c}_\phi^\dagger \mathbf{c}_\psi).
\end{equation}

\begin{figure}[t]
\centering
\includegraphics[width=0.8\linewidth]{Interference_Based_decision_Geometry.png}
\caption{Comparison of decision boundaries: Linear Interference (left) vs. Fidelity-based Quadratic (right). The interference rule $\mathrm{Re}\langle \phi | \psi \rangle = 0$ induces a linear hyperplane and retains the sign of the real interference term, while fidelity-based rules result in curved boundaries that discard sign information.}
\label{fig:geometry_comparison}
\end{figure}

\begin{figure}[H]
\centering
\includegraphics[width=0.8\linewidth]{Fidelity_based_similarity_geometry.png}
\caption{Detailed geometric profile of fidelity-based similarity measures in Hilbert space, illustrating the nonlinear curvature of the decision surface.}
\label{fig:fidelity_geometry}
\end{figure}

\subsubsection{Formal Derivation of the Decision Score}
Consider the normalized state $|\psi\rangle = \sum_i a_i |i\rangle$ and the representative state $|\phi\rangle = \sum_i b_i |i\rangle$. The interference term $\text{Re}\langle \phi | \psi \rangle$ can be expanded in terms of the real and imaginary parts of the coefficients:
\begin{equation}
    \langle \phi | \psi \rangle = \sum_i b_i^* a_i = \sum_i (\text{Re } b_i - i \text{Im } b_i)(\text{Re } a_i + i \text{Im } a_i)
\end{equation}
Taking the real part, we obtain:
\begin{equation}
    \text{Re}\langle \phi | \psi \rangle = \sum_i (\text{Re } a_i \text{Re } b_i + \text{Im } a_i \text{Im } b_i)
\end{equation}
This summation is identical to the Euclidean inner product in $\mathbb{R}^{2 \cdot 2^n}$. The decision boundary $\text{Re}\langle \phi | \psi \rangle = 0$ is therefore the locus of all states in $\mathcal{H}$ that are "mutually orthogonal" under this real-part inner product.

\subsubsection{Equivalence to a Classical Linear Classifier}
\label{sec:equivalence}
It is important to state precisely what the interference primitive computes, so that its near-term advantage is not overstated. In our pipeline the CNN neck terminates in a ReLU activation, so every embedding---and hence every amplitude-encoded state $|\psi\rangle$---is \emph{real} and non-negative. For real $\phi,\psi$ the imaginary machinery of the overlap vanishes and $\mathrm{Re}\langle\phi|\psi\rangle = \sum_i \phi_i\psi_i = \bm{\phi}\cdot\bm{\psi}$: the decision statistic is exactly the Euclidean dot product. Consequently:
\begin{itemize}
    \item The static IQC of Section~\ref{sec:centroids}, with $\bm{\chi}\propto\bm{\phi}^{(+)}-\bm{\phi}^{(-)}$, is exactly the signed-centroid rule $\hat{y}=\mathrm{sign}(\bm{\chi}\cdot\bm{x})$---a single hyperplane whose normal is the difference of class centroids, in the spirit of linear discriminant analysis with identity covariance.
    \item A multi-prototype model that aggregates memory overlaps by a weighted vote, $\sum_i w_i(\bm{m}_i\cdot\bm{x}) = \big(\sum_i w_i\bm{m}_i\big)\cdot\bm{x}$, collapses at inference to a \emph{single} effective hyperplane $\bm{W}=\sum_i w_i\bm{m}_i$, despite storing $2K$ prototype states.
\end{itemize}
We verify this empirically: the statevector Hadamard-test circuit reproduces the numpy dot product to $\max|\Delta|=1.2\times10^{-12}$, a high-shot circuit agrees within shot noise, and both the static rule and the frozen $K{=}3$ model match their linear replicas on $100\%$ of test predictions. The practical implication is that the IQC cannot exceed a well-regularized classical linear classifier on the same features; its genuine near-term contribution is therefore the \emph{measurement-efficient sign readout} (Section~\ref{sec:shots}) and the geometric flexibility introduced by \emph{multiple} prototypes (Section~\ref{sec:capacity}), not a supra-classical decision function.

\subsection{Quantum Interference Circuit}
The accessibility of the linear overlap signal $\mathrm{Re}\langle \phi | \psi \rangle$ is the defining hardware requirement of our framework. Since standard projective measurements $\langle \psi | O | \psi \rangle$ are quadratic in the state amplitudes, we utilize an ancilla-assisted interference circuit. In this framework, one inference corresponds to executing a fixed interference circuit on an input state and assigning a label based on the sign of the estimated overlap.

\subsubsection{Physical Realization via the Hadamard Test}
As with all gate-based quantum classifiers, state-preparation is treated as a separate cost; the inference complexity discussed here refers exclusively to the decision observable and measurement process.

The real part of the overlap between two quantum states can be directly extracted using the Hadamard test \cite{ekert2002direct}. The circuit requires one additional ancilla qubit and the original register used to store the data states.

The sequence for extracting $\mathrm{Re}\langle \phi | \psi \rangle$ via the Hadamard test is summarized below \cite{ekert2002direct}:
\begin{enumerate}
    \item \textbf{State Initialization}: Both the ancilla and the data register are reset; the register is then loaded with the input state $|\psi\rangle$.
    \item \textbf{Superposition Creation}: The ancilla is placed into an equal superposition $\frac{1}{\sqrt{2}}(|0\rangle + |1\rangle)$ by a Hadamard transformation.
    \item \textbf{Controlled Transition}: A controlled unitary acts on the register, conditioned on the ancilla, effecting the map $|\psi\rangle \mapsto |\phi\rangle$. When both states are generated from $|0\rangle$ by unitaries $V_\psi$ and $V_\phi$, this reduces to the controlled operator $V_\phi V_\psi^\dagger$.
    \item \textbf{Interference Gate}: A second Hadamard is applied to the ancilla to convert the relative phase into a measurable amplitude difference.
    \item \textbf{Ancilla Readout}: The ancilla qubit is measured in the computational basis, yielding the interference signal.
\end{enumerate}

The probability of measuring the ancilla in state $|0\rangle$ is:
\begin{equation}
P(0) = \frac{1}{2} + \frac{1}{2} \mathrm{Re}\langle \phi | \psi \rangle.
\end{equation}
The interference signal (our decision score) is then obtained from the expectation value of the ancilla's $Z$-observable:
\begin{equation}
\langle Z_{\text{anc}} \rangle = P(0) - P(1) = \mathrm{Re}\langle \phi | \psi \rangle.
\end{equation}

\begin{figure}[t]
\centering
\includegraphics[width=1.0\linewidth]{transition_isdo.png}
\caption{Circuit diagram of the ISDO-B implementation. The controlled transition unitary $U_{\chi\psi}$ is used to estimate the linear overlap $\mathrm{Re}\langle \phi | \psi \rangle$ by measuring the ancilla qubit expectation value over a small constant number of shots.}
\label{fig:circuit}
\end{figure}

\subsubsection{ISDO-B: Hardware-Efficient Linear Interference}
While the canonical Hadamard test relies on controlled state-preparation for both $|\psi\rangle$ and $|\phi\rangle$, the \textbf{ISDO-B (Interference-Based Static Decision Observable)} variant utilizes a controlled transition unitary $U_{\chi\psi}$ where $U_{\chi\psi}|\psi\rangle = |\phi\rangle$. We assume $U_{\chi\psi}$ is a state-transfer operator acting such that the target subspace is correctly mapped; its action on the orthogonal subspace is unconstrained as it does not contribute to the interference signal.

The ISDO-B circuit adapts the Hadamard test as follows:
\begin{enumerate}
    \item Both ancilla and data register are initialized to $|0\rangle$; the register is then set to the test state $|\psi\rangle$.
    \item A Hadamard transformation brings the ancilla into superposition, producing the joint state $\frac{1}{\sqrt{2}}(|0\rangle+|1\rangle) \otimes |\psi\rangle$.
    \item The controlled transition unitary $CU_{\chi\psi}$ is engaged: when the ancilla is $|1\rangle$, the register evolves from $|\psi\rangle$ to $|\phi\rangle$, yielding $\frac{1}{\sqrt{2}}(|0\rangle |\psi\rangle + |1\rangle |\phi\rangle)$.
    \item A final Hadamard on the ancilla maps the overlap phase onto the $Z$-basis, after which the ancilla is measured.
\end{enumerate}
The expectation value of the ancilla measurement yields $\langle Z_{\text{anc}} \rangle = \mathrm{Re}\langle \phi | \psi \rangle$. This formulation is particularly advantageous in static regimes where the class prototype $|\phi\rangle$ is precomputed, allowing for the engineering of specialized observables that avoid the overhead of full controlled state-preparation of arbitrary centroids.

\subsubsection{Single-Circuit Inference Cost}
The IQC uses a single fixed interference circuit per sample and reads out one ancilla qubit to obtain the decision sign. In our experiments, a small constant number of shots is sufficient to stabilize the sign decision, in contrast to fidelity-based methods that require high-precision probability estimation.

\subsection{Class Representative State Construction}
\label{sec:centroids}
A pivotal design choice in the IQC framework is the complete decoupling of the learning process from quantum execution. Unlike VQCs, where learning involves an iterative quantum-classical loop, our framework performs all structural learning classically before anything is loaded onto the quantum processor.

\subsubsection{Centroid-Based State Construction}
The learning objective is to find a normalized quantum state $|\phi^{(c)}\rangle$ that best represents the distribution of training samples belonging to class $c \in \{-1, +1\}$. In this work, we adopt a \textbf{centroid-based aggregation rule}, which is conceptually equivalent to calculating the mean vector in Hilbert space. This centroid-based approach is conceptually similar to classical linear discriminant analysis methods, where class means are used to define approximate linear decision boundaries, albeit implemented directly within the high-dimensional quantum Hilbert space.

Given a subset of training states $\mathcal{S}_c = \{|\psi_i\rangle : y_i = c\}$, where each $|\psi_i\rangle$ is an amplitude-encoded vector in $\mathbb{C}^d$, the unnormalized class-representative vector is:
\begin{equation}
| \tilde{\phi}^{(c)} \rangle = \sum_{|\psi_i\rangle \in \mathcal{S}_c} w_i |\psi_i\rangle,
\end{equation}
where $w_i$ are importance weights. In the simplest case, where all training samples are considered equally, $w_i = 1/|\mathcal{S}_c|$. The finalized class-representative state is obtained by normalization:
\begin{equation}
|\phi^{(c)}\rangle = \frac{| \tilde{\phi}^{(c)} \rangle}{\| | \tilde{\phi}^{(c)} \rangle \|}.
\end{equation}

Algorithm~\ref{alg:learning} details the classical preprocessing steps required to construct these states.

\begin{algorithm}[ht]
\caption{Static Learning of Class-Representative States}
\label{alg:learning}
\begin{algorithmic}[1]
\REQUIRE Training dataset $\mathcal{D} = \{(\mathbf{x}_i, y_i)\}_{i=1}^N$
\STATE Select feature map $\Phi: \mathbb{R}^d \rightarrow \mathbb{C}^{2^n}$ (e.g., Amplitude Encoding)
\STATE Initialize class accumulators: $\mathbf{v}_+ \leftarrow \mathbf{0}, \mathbf{v}_- \leftarrow \mathbf{0}$
\FOR{each $(\mathbf{x}_i, y_i) \in \mathcal{D}$}
    \STATE Compute state vector: $|\psi_i\rangle \leftarrow \Phi(\mathbf{x}_i)$
    \IF{$y_i = +1$}
        \STATE $\mathbf{v}_+ \leftarrow \mathbf{v}_+ + |\psi_i\rangle$
    \ELSE
        \STATE $\mathbf{v}_- \leftarrow \mathbf{v}_- + |\psi_i\rangle$
    \ENDIF
\ENDFOR
\STATE Normalize: $|\phi^{(+)}\rangle \leftarrow \mathbf{v}_+/\|\mathbf{v}_+\|, |\phi^{(-)}\rangle \leftarrow \mathbf{v}_-/\|\mathbf{v}_-\|$
\RETURN $|\phi^{(+)}\rangle, |\phi^{(-)}\rangle$
\end{algorithmic}
\end{algorithm}

This aggregation is performed entirely in classical memory using the complex amplitude vectors of the data embeddings. For a dataset of size $N$ and a feature space of dimension $d$, this classical step scales as $O(Nd)$, which is significantly faster than the $O(M \cdot N)$ scaling of VQC training (where $M$ is the number of optimization iterations).

\subsubsection{State Selection and Prototype Extraction}
In some applications, a simple average may not be the most representative summary of a class, especially if the class distribution is multimodal. However, within the scope of this work, we assume a single-mode distribution for each class. The state $|\phi^{(c)}\rangle$ essentially acts as a "prototype" or "centroid" in the quantum feature space. The linear interference score $\text{Re}\langle \phi | \psi \rangle$ can then be interpreted as the projection of the test sample $|\psi\rangle$ onto this class prototype.

\subsubsection{Static Nature and Hybrid Efficiency}
The resulting states $|\phi^{(+)}\rangle$ and $|\phi^{(-)}\rangle$ are static. Once computed, they are used to derive the control sequences for the interference circuit described earlier. There are no weight updates, no memory growth, and no backpropagation through quantum gates.

By treating the quantum device as a static inference engine rather than a trainable model, we eliminate several pathologies common in hybrid QML:
\begin{itemize}
    \item \textbf{Gradient Noise}: Classical aggregation is deterministic and does not suffer from the sampling variance of quantum-derived gradients.
    \item \textbf{Barren Plateaus}: We avoid the optimization landscape altogether. The "training" is a simple linear aggregation with a guaranteed closed-form solution.
    \item \textbf{I/O Overhead}: The classical-quantum bridge is traversed only for data loading and the final inference shot, avoiding the high-latency loops of variational training.
\end{itemize}

\section{Implementation}

\subsection{Data Source: PatchCamelyon}
For our experiments, we utilize the PatchCamelyon (PCam) dataset, a cornerstone for benchmarking automated detection of lymph node metastases \cite{pcam_veeling2018}. PCam comprises 96\,$\times$\,96 RGB histopathology patches, where each patch is labeled by whether metastatic tumor cells appear in the central 32\,$\times$\,32 pixel region. This binary classification task differentiates between tumor-infiltrated (malignant) and non-tumor (benign) tissue. The dataset's complexity stems from the intricate morphologies of lymphatic structures, where malignant cells disrupt normalized tissue architecture. Representative examples of these tissue types are shown in Fig.~\ref{fig:benign_malignant}.

\begin{figure}[ht]
    \centering
    \includegraphics[width=0.45\linewidth]{benign.png}
    \includegraphics[width=0.45\linewidth]{malignant.png}
    \caption{Representative lymph node histopathology patches from the PCam dataset. Left: Benign tissue showing normal lymphatic architecture. Right: Malignant tissue exhibiting disrupted cellular morphology and tumor infiltration.}
    \label{fig:benign_malignant}
\end{figure}

\subsection{CNN Feature Extraction}
The CNN backbone and neck are trained classically using cross-entropy loss on PCam. We follow the original PCam train/val/test split provided by the dataset with strict separation of roles: the backbone is trained on the \texttt{train} split (262{,}144 patches), the \texttt{val} split (32{,}768 patches) is used only for model selection, and \textbf{all final downstream numbers are reported on the held-out \texttt{test} split} (32{,}768 patches), which is touched only at evaluation time. This corrects an earlier protocol that both selected and evaluated on a subset of the validation partition; on the true test split every method's accuracy is lower, reflecting the well-known PCam validation-to-test distribution shift. To support statistical validation, we retrain the CNN from scratch under five random seeds and report mean\,$\pm$\,std across seeds for all downstream comparisons.

Training uses a lightweight PCamCNN with three convolutional blocks (3$\times$3 conv + batch normalization + ReLU, with max-pooling after the first two blocks), followed by adaptive average pooling and a linear embedding head of dimension $d=32$ with ReLU activation, inspired by prior hybrid quantum imaging pipelines \cite{pcam_01,pcam_03}. A temporary linear classifier is attached during training and discarded after convergence. We use Adam (learning rate $10^{-3}$, weight decay $10^{-4}$), batch size 64, and a ReduceLROnPlateau scheduler (factor 0.5, patience 2) with early stopping (patience 10). Data augmentation is limited to label-preserving transforms: random horizontal/vertical flips and light color jitter, followed by normalization to mean $[0.5,0.5,0.5]$ and std $[0.5,0.5,0.5]$. Validation and embedding extraction use deterministic normalization only.

After training, the CNN+neck (PCamCNN) is frozen and used to produce embeddings for all training samples. These embeddings are then aggregated into class prototypes (Section~\ref{sec:centroids}). The quantum processor is used only at inference.

\subsection{Quantum Encoding Setup}
In our experiments, $d = 32$ and $n = 5$, so the quantum register consists of 5 data qubits and 1 ancilla qubit (6 total). Each embedding $\mathbf{x} \in \mathbb{R}^{32}$ is amplitude-encoded using Qiskit's \texttt{StatePreparation} unitary, which realizes the mapping $|\psi\rangle = \sum_{i=0}^{31} x_i|i\rangle$. The \texttt{StatePreparation} algorithm prepares arbitrary normalized complex vectors via uniformly controlled rotations \cite{plesch2011quantum}.

\subsubsection{Depth Analysis of Amplitude Encoding}
While the IQC inference circuit is fixed and short-depth (constant relative to state-preparation), the overall performance is bounded by the cost of preparing the states $|\psi\rangle$ and $|\phi\rangle$. For a $d$-dimensional input vector, amplitude encoding into $n = \log_2 d$ qubits requires a unitary $U$ such that $U|0\rangle^{\otimes n} = \sum_{i=1}^d x_i |i\rangle$.

The standard decomposition of such a unitary using the Plesch-Brukner algorithm results in a circuit with $O(d)$ CNOT gates and $O(d)$ single-qubit rotations. Its asymptotic complexity was rigorously analyzed by Plesch and Brukner \cite{plesch2011quantum}. While this is exponential in the number of qubits $n$, it is linear in the number of features $d$. For NISQ deployment, one often uses approximate state-preparation techniques or hardware-efficient ans\"atze that trade off exact representation for reduced depth. In our simulation, we consider an exact decomposition but point out that the linear interference mechanism of the IQC is still valid independently of the state-preparation technique used, as long as the relative phase information is maintained.

\subsection{Simulation Environment}
Simulations are carried out via a statevector quantum simulator, with a physically wired \texttt{qiskit-aer} noise model used for the noise-robustness study. The shot complexity is measured relative to the physical measurement step, not simulation time. All downstream results are reported as mean\,$\pm$\,std over five random seeds, where each seed controls CNN initialization, data sub-sampling, prototype clustering, online update order, and shot sampling through independent sub-streams derived from a single seed sequence.

All simulations were implemented using a Python-based stack, where Qiskit was used for circuit design and statevector simulation, and PyTorch for the classical CNN modules. Classical training and embedding extraction were performed on an NVIDIA RTX PRO 6000 GPU; quantum circuit execution used the CPU \texttt{AerSimulator}, which is optimal at the six-qubit scale of the ISDO circuit (GPU statevector simulation only pays off at substantially larger qubit counts). The simulation environment details are as follows:
\begin{itemize}
    \item \textbf{Data Handling}: Downstream models are fit on a balanced sub-sample of the \texttt{train} split and evaluated on the full 32{,}768-patch \texttt{test} split; the fidelity-kernel and variational baselines are capped at a smaller fit/evaluation size for tractability, as noted in their results.
    \item \textbf{Dimensionality Reduction}: The CNN backbone (PCamCNN) compresses each RGB patch into a $d=32$ embedding, which is L2-normalized prior to quantum conversion.
    \item \textbf{Quantum Registry}: Embeddings are amplitude-encoded into a 6-qubit system ($n=5$ data qubits and 1 ancilla) via Qiskit's \texttt{StatePreparation} routine.
    \item \textbf{Noise Simulation}: We apply a depolarizing model (single- and two-qubit depolarizing errors attached to the transpiled native basis) and a thermal-relaxation ($T_1, T_2$) model, both injected through \texttt{qiskit-aer}'s \texttt{NoiseModel} and transpiled against the noisy basis so that errors attach to the decomposed \texttt{StatePreparation} gates. An analytic tier applies the corresponding $(1-p)^L$ attenuation to the exact overlap for full-test-set coverage, validated against the circuit tier on a subset.
\end{itemize}

\subsection{Benchmark Models}
Our IQC (ISDO-B variant) is compared against classical baselines (Logistic Regression, Linear SVM, $k$-NN) and measurement-based quantum baselines (a fidelity quantum-kernel SVM, QSVM; and a variational quantum classifier, VQC) on the same CNN embeddings, together with the multi-prototype (fixed and adaptive) interference regimes. All methods consume identical per-seed embeddings and are evaluated on the full \texttt{test} split, except QSVM and VQC, which are capped at 3{,}500 training and 2{,}000 evaluation samples for tractability.

\subsection{Main Comparison}
\label{sec:main}
Table~\ref{tab:metrics} reports classification accuracy on the held-out \texttt{test} split as mean\,$\pm$\,std over five seeds, each with an independently retrained CNN. Statistical significance of pairwise differences is assessed with a per-seed paired McNemar test on the $32{,}768$ shared test predictions; with this sample size every pairwise difference is significant except Linear SVM vs.\ Logistic Regression ($p=0.53$).

\begin{table}[ht]
\centering
\caption{Classification performance on the held-out PCam \texttt{test} split (five seeds, mean\,$\pm$\,std). QSVM/VQC are capped at 3{,}500/2{,}000 fit/eval samples.}
\label{tab:metrics}
\begin{tabular}{l c}
\toprule
Method & Accuracy (\%) \\
\midrule
Linear SVM & $84.05 \pm 1.12$ \\
Logistic Regression & $84.04 \pm 1.15$ \\
$k$-NN ($k=5$) & $83.36 \pm 1.14$ \\
QSVM (fidelity kernel) & $83.21 \pm 0.97$ \\
\textbf{IQC static (Proposed)} & $\mathbf{82.67 \pm 1.35}$ \\
IQC fixed ($K{=}3$) & $81.36 \pm 1.96$ \\
IQC adaptive & $80.89 \pm 2.80$ \\
VQC & $71.62 \pm 2.66$ \\
\bottomrule
\end{tabular}
\end{table}

On honest test data the single-prototype IQC ($82.67\%$) is the strongest member of the interference family, decisively outperforms the variational baseline, and is statistically comparable to the fidelity quantum kernel, while trailing the well-tuned classical linear methods by $\sim\!1.4$ points. This gap is expected: by the equivalence of Section~\ref{sec:equivalence}, the IQC \emph{is} a linear classifier and therefore cannot exceed an optimally regularized linear classifier on the same features. Its value proposition is measurement efficiency (Section~\ref{sec:shots}) and the multi-prototype geometry of Section~\ref{sec:capacity}.

\subsection{Shot-Complexity Behavior}
\label{sec:shots}
The measurement-efficiency claim is confirmed and, importantly, is governed by the empirical decision margin rather than the model's memory count. Because the ancilla readout is a Bernoulli variable with $P(0)=\tfrac12(1+s)$ for overlap $s$, the shots required to fix the sign with confidence $1-\delta$ obey the Hoeffding bound $N \gtrsim 2\ln(1/\delta)/m^2$ at margin $m=|s|$---independent of the feature dimension. The single-prototype static IQC has a large median margin ($m\approx0.34$) and reaches its exact-backend accuracy within Monte-Carlo error by $\sim\!64$--$128$ shots; the $K{=}3$ model has a smaller median margin ($m\approx0.13$) and therefore needs more shots for the same stability. Empirical sampled accuracy lies between the exact binomial sign-flip prediction and the (looser) Hoeffding bound, converging to the noiseless accuracy at high shot counts, as shown in Fig.~\ref{fig:shots}.

\begin{figure}[t]
\centering
\includegraphics[width=0.85\linewidth]{shot_sweep_iqc_static.png}
\caption{Accuracy of the single-prototype IQC versus measurement shots (log scale) on the full test split, with the exact binomial prediction and the exact-backend accuracy overlaid. Near-peak accuracy is reached at a small constant shot count.}
\label{fig:shots}
\end{figure}

\subsection{Prototype Capacity and Adaptive Regimes}
\label{sec:capacity}
Section~\ref{sec:centroids} noted that a single centroid may poorly summarize a multimodal class. We test this directly by sweeping the number of prototypes per class, $K$, in a fixed-memory interference model (winner-take-all training over $K$-means prototypes) and comparing against $k$-NN on the same features. As shown in Table~\ref{tab:capacity} and Fig.~\ref{fig:capacity}, test accuracy grows monotonically with $K$ and \textbf{surpasses $k$-NN at $K\!\geq\!11$} (paired McNemar $p<10^{-25}$), rising from $80.2\%$ at $K{=}1$ to $83.9\%$ at $K{=}23$. This confirms that the single-mode centroid, not the interference measurement, is the accuracy bottleneck, and that additional prototypes---all still evaluated by the same fixed sign-readout circuit---recover the discriminative structure. An adaptive regime that spawns and prunes prototypes online reaches comparable accuracy with fewer than half the memories of a fixed $K{=}3$ model, indicating that the number of prototypes can be learned rather than fixed.

\begin{table}[ht]
\centering
\caption{Prototype-capacity sweep (test split, five seeds, mean accuracy \%). $k$-NN($5$) reference $=83.4$.}
\label{tab:capacity}
\begin{tabular}{l c c c c c c}
\toprule
$K$ & 1 & 3 & 5 & 7 & 11 & 23 \\
\midrule
IQC & 80.2 & 80.6 & 81.5 & 81.8 & 83.5 & 83.9 \\
\bottomrule
\end{tabular}
\end{table}

\begin{figure}[t]
\centering
\includegraphics[width=0.85\linewidth]{capacity_sweep.png}
\caption{Test accuracy of the multi-prototype IQC versus prototypes-per-class $K$ (mean\,$\pm$\,std over five seeds), with the $k$-NN reference. The interference classifier crosses the $k$-NN line at $K\!\geq\!11$.}
\label{fig:capacity}
\end{figure}

\section{Discussion}

\subsection{Measurement Efficiency}
The IQC reaches near-peak performance with small constant shot counts, while fidelity-based and VQC models require substantially more measurements to stabilize classification scores. This reduces the duty cycle pressure on quantum hardware and improves throughput in measurement-limited settings.

The reduction in measurement shots has a direct impact on the energy consumption of the quantum-classical system. In superconducting architectures, each measurement shot involves microwave pulses and high-speed classical digitization, consuming significant power at the control electronics layer. By reducing the number of shots from $10^4$ (typical for fidelity estimation) to a small constant number of shots (often $10^1-10^2$ for IQC sign estimation), we achieve a reduction in the "energy per inference" metric. This potentially positions interference-based QML as a more energy-efficient alternative for data processing in future quantum deployments.

\subsection{Noise Resilience Analysis}
In typical NISQ environments, circuits are prone to operational errors and decoherence. We evaluate the IQC's performance under two standard noise channels.

\subsubsection{Depolarizing Effects}
The depolarizing channel is modeled as $\mathcal{E}(\rho) = (1-p)\rho + p\frac{I}{2^{n+1}}$, i.e., the state is replaced by the maximally mixed state with probability $p$ \cite{nielsen2000qcqi}. For our interference circuit, this channel attenuates the overlap signal by a factor of $(1-p)^L$, where $L$ denotes circuit depth:
\begin{equation}
    \langle Z_{\text{anc}} \rangle_{\text{noisy}} = (1-p)^L \mathrm{Re}\langle \phi | \psi \rangle.
\end{equation}
Notably, since the attenuation factor $(1-p)^L$ remains positive, the \textbf{sign} of the expectation value is preserved \emph{in expectation}. However, this sign-preservation argument must be read together with the actual value of $L$. Exact amplitude encoding of an arbitrary $32$-dimensional state, together with the controlled state-preparations of the Hadamard test, transpiles (native basis $\{\mathrm{sx},\mathrm{x},\mathrm{rz},\mathrm{cx}\}$) to a circuit of depth $L\approx1268$, comprising roughly $1128$ single-qubit and $462$ two-qubit gates---the cost is dominated by the controlled state-preparation, not the ancilla. At this depth the multiplicative attenuation is severe: with a two-qubit error rate of $10p$, the overlap is reduced to $\sim\!56\%$ of its value at $p=10^{-4}$ and to $\sim\!0.3\%$ at $p=10^{-3}$, at which point the sign is no longer recoverable from a bounded shot budget and accuracy falls to chance. Table~\ref{tab:noise} and Fig.~\ref{fig:noise} report the resulting \textbf{survival threshold near $p\approx10^{-4}$}; the analytic $(1-p)^L$ tier and the \texttt{qiskit-aer} circuit tier agree (both at chance for $p\geq10^{-3}$). We therefore refine the earlier resilience claim: the sign statistic is intrinsically robust, but for exact-encoding circuits the operative depth $L\sim10^{3}$ places the usable depolarizing regime at $p\lesssim10^{-4}$, well below the $p=10^{-2}$ operating point of many NISQ devices. Realizing the resilience the primitive theoretically enjoys thus requires shallow or approximate state-preparation rather than exact decomposition. Thermal relaxation shows the same qualitative picture: signal is lost near $T_1\approx30\,\mu$s and recovers only for $T_1\gtrsim100$--$200\,\mu$s.

\begin{table}[ht]
\centering
\caption{Depolarizing survival threshold for the single-prototype IQC (test split, $1024$ shots). Analytic and circuit tiers agree.}
\label{tab:noise}
\begin{tabular}{l c c c c c}
\toprule
$p$ & $10^{-5}$ & $10^{-4}$ & $3\times10^{-4}$ & $10^{-3}$ & $10^{-2}$ \\
\midrule
Attenuation & 0.94 & 0.56 & 0.18 & 0.003 & $\sim\!0$ \\
Accuracy & 0.824 & 0.824 & 0.797 & 0.508 & 0.499 \\
\bottomrule
\end{tabular}
\end{table}

\begin{figure}[t]
\centering
\includegraphics[width=0.85\linewidth]{noise_depolarizing_static.png}
\caption{Single-prototype IQC accuracy versus depolarizing rate $p$ (log scale) on the test split. The signal survives only up to $p\approx10^{-4}$ at the exact-encoding circuit depth $L\approx1268$; the analytic and \texttt{qiskit-aer} circuit tiers agree.}
\label{fig:noise}
\end{figure}

\subsubsection{Dephasing Errors}
Dephasing or phase-damping errors, $\mathcal{E}_{ph}(\rho) = (1-p)\rho + p Z\rho Z$, may arise during unitary execution or idle periods \cite{nielsen2000qcqi}. Within the context of the Hadamard test, dephasing on the ancilla adds a stochastic phase $\theta$ to the interference term:
\begin{equation}
    P(0) = \frac{1}{2} + \frac{1}{2} \text{Re}(e^{i\theta} \langle \phi | \psi \rangle).
\end{equation}
Unbiased dephasing noise typically results in a stable expected decision score; however, systematic phase drifts can shift the decision boundary, suggesting that the ISDO-B variant is more resilient to stochastic depolarization than to fixed calibration errors.

\subsection{Limitations of the Model}
A fundamental limitation of the current IQC framework is its reliance on a linear decision boundary in Hilbert space. While the Hilbert space dimension $2^n$ is large, separation depends heavily on the quality of the feature map $\Phi(\mathbf{x})$. If the data distribution is not linearly separable under the chosen map, the IQC behaves like a high-bias linear classifier. The dependence of inference stability on the classification margin mirrors that of classical linear classifiers and reflects intrinsic ambiguity near the decision boundary rather than a limitation of the interference-based inference mechanism.

A second, more fundamental limitation follows from Section~\ref{sec:equivalence}: because the decision rule reduces to a linear classifier in the embedding space, the IQC inherits the ceiling of linear separability under the chosen feature map and cannot, on accuracy alone, improve upon a well-regularized classical linear model on the same embeddings. Its advantage is confined to measurement efficiency and to the geometric capacity added by multiple prototypes.

All experiments reported in this work are conducted entirely via classical statevector simulation. No real quantum hardware experiments have been performed. Noise is injected via \texttt{qiskit-aer} depolarizing and thermal-relaxation models transpiled against the native basis; while this attaches errors to the decomposed circuit faithfully, it still reflects idealized channels rather than the full complexity of real device errors (e.g., crosstalk, leakage, coherent errors). Relative to the earlier version of this study, we have addressed two previously identified weaknesses---single-trial reporting and evaluation on a validation subset---by adopting five-seed mean\,$\pm$\,std reporting with paired significance testing and by evaluating exclusively on the held-out \texttt{test} split. Real hardware validation on a backend such as IBM Quantum or IonQ, and shallow/approximate state-preparation to bring the operative circuit depth into the noise-tolerant regime, remain the most critical directions for strengthening and extending these claims.

\subsection{Scalability Considerations}
The transition of the IQC from simulation to actual quantum hardware needs to take into account the issue of gate decomposition and ancilla usage. The Hadamard test needs controlled-unitary operations $CU$, which can be expensive in terms of circuit depth $L$ when decomposed into native gates. The interference circuit relies on a single ancilla qubit that must be connected to all qubits in the data register used in the controlled operations; restricted connectivity can introduce SWAP overhead. Readout error rates can be as high as 1--5\% in current NISQ devices, but the IQC's reliance on a single ancilla measurement helps isolate this effect. Readout error mitigation techniques can be applied to the ancilla qubit to improve the precision of the interference score without multi-qubit tomography.

The IQC framework exhibits favorable scaling compared to both classical linear models and variational quantum models. To represent a $d$-dimensional feature vector, the IQC requires $n = \lceil \log_2 d \rceil$ qubits, achieving logarithmic qubit scaling in representation (with linear gate complexity for state-preparation). Once the state is prepared, the interference circuit depth $L$ is a constant or near-constant factor relative to the state-preparation depth. Inference uses a single fixed circuit and a small constant number of shots to stabilize the sign decision. The classical aggregation step scales as $O(Nd)$, which is optimal for linear learning. The reduction in measurement shots also lowers the energy per inference in measurement-limited hardware settings.

\section{Concluding Remarks}
In this research, we have proposed and rigorously characterized an interference-driven quantum classifier designed as a static inference primitive for histopathological image assessment. By centering the decision rule on the linear overlap $\mathrm{Re}\langle \phi | \psi \rangle$, we obtain a measurement-efficient alternative to quadratic methods whose shot budget is governed by the decision margin. Our analysis clarifies the nature of the primitive: under real amplitude encoding it is exactly a linear (signed-centroid) classifier, so its near-term value lies in the measurement-efficient sign readout and in the geometric capacity of multiple prototypes rather than in a supra-classical decision function. On the held-out PCam test split with five-seed validation, the single-prototype IQC ($82.7\pm1.4\%$) is competitive with a fidelity quantum kernel and trails well-tuned classical linear baselines by a small, statistically significant margin, while a multi-prototype extension surpasses $k$-NN for $K\!\geq\!11$. We also give a corrected account of noise robustness: the sign statistic is intrinsically stable, but the exact-encoding circuit depth ($L\approx1268$) confines usable depolarizing noise to $p\lesssim10^{-4}$, so realizing the primitive's theoretical resilience requires shallow or approximate state-preparation.

Subsequent research will pursue three directions suggested by these findings: shallow and approximate state-preparation to move the operative circuit depth into the noise-tolerant regime; learned prototype budgets via the adaptive spawn/prune regime, which already matches fixed-memory accuracy with fewer prototypes; and hardware validation on superconducting or trapped-ion backends where the single-ancilla sign readout can be measured directly.

\subsection{Extension to Multi-class Classification}
While this work focuses on binary classification as the core primitive, the IQC framework can be extended to multi-class problems ($K > 2$) using standard decomposition strategies such as \textbf{One-vs-All (OvA)} or \textbf{One-vs-One (OvO)}.

In the OvA approach, for $K$ classes, we learn $K$ static class-representative states $\{|\phi^{(1)}\rangle, |\phi^{(2)}\rangle, \dots, |\phi^{(K)}\rangle\}$. The interference operation is performed between the test state $|\psi\rangle$ and each class prototype. The classification rule is then:
\begin{equation}
    \hat{y} = \arg\max_c \mathrm{Re}\langle \phi^{(c)} | \psi \rangle.
\end{equation}
The shot complexity for multi-class inference scales as $O(K)$, as we must perform $K$ independent interference measurements. However, each measurement uses a small number of shots, so the overall efficiency remains significantly higher than multi-class VQCs or QSVMs. Future work will investigate more sophisticated ways to load multiple class prototypes into a single quantum memory for parallel interference, potentially reducing the scaling to $O(\log K)$.

\clearpage
\bibliographystyle{IEEEtran}
\begin{thebibliography}{99}

\bibitem{benedetti2019parameterized}
Marcello Benedetti, Erika Lloyd, Stefan Sack, and Mattia Fiorentini, ``Parameterized quantum circuits as machine learning models,'' \emph{Quantum Science and Technology}, vol.~4, no.~4, Art. no.~043001, 2019, doi: 10.1088/2058-9565/ab0536.

\bibitem{mcclean2018barren}
Jarrod R. McClean, Sergio Boixo, Vadim N. Smelyanskiy, Ryan Babbush, and Hartmut Neven, ``Barren plateaus in quantum neural network training landscapes,'' \emph{Nature Communications}, vol.~9, Art. no.~4812, 2018, doi: 10.1038/s41467-018-07090-4.

\bibitem{schuld2019feature}
Maria Schuld and Nathan Killoran, ``Quantum machine learning in feature Hilbert spaces,'' \emph{Physical Review Letters}, vol.~122, no.~4, Art. no.~040504, 2019, doi: 10.1103/PhysRevLett.122.040504.

\bibitem{schuld2017distance}
Maria Schuld, Mark Fingerhuth, and Francesco Petruccione, ``Implementing a distance-based classifier with a quantum interference circuit,'' \emph{EPL (Europhysics Letters)}, vol.~119, no.~6, Art. no.~60002, 2017, doi: 10.1209/0295-5075/119/60002.

\bibitem{abbas2021power}
Amira Abbas, David Sutter, Christa Zoufal, Aur\'elien Lucchi, Alessio Figalli, and Stefan W\"{o}rner, ``The power of quantum neural networks,'' \emph{Nature Computational Science}, vol.~1, no.~6, pp.~403--409, 2021, doi: 10.1038/s43588-021-00084-1.

\bibitem{havlicek2019supervised}
Vojt\v{e}ch Havl\'{\i}\v{c}ek, Antonio D. C\'orcoles, Kristan Temme, Aram W. Harrow, Abhinav Kandala, Jerry M. Chow, and Jay M. Gambetta, ``Supervised learning with quantum-enhanced feature spaces,'' \emph{Nature}, vol.~567, no.~7747, pp.~209--212, 2019, doi: 10.1038/s41586-019-0980-2.

\bibitem{sergioli2019binary}
Giuseppe Sergioli, Roberto Giuntini, and H\'ector Freytes, ``A new quantum approach to binary classification,'' \emph{PLoS ONE}, vol.~14, no.~5, e0216224, 2019, doi: 10.1371/journal.pone.0216224.

\bibitem{gonzalez2021classification}
Fabio A. Gonz\'alez, Vladimir Vargas-Calder\'on, and Herbert Vinck-Posada, ``Classification with quantum measurements,'' arXiv:2004.01227v2, 2021.

\bibitem{blank2020kernel}
Carsten Blank, Daniel K. Park, June-Koo Kevin Rhee, and Francesco Petruccione, ``Quantum classifier with tailored quantum kernel,'' \emph{npj Quantum Information}, vol.~6, Art. no.~41, 2020, doi: 10.1038/s41534-020-0272-6.

\bibitem{blank2022compact}
Carsten Blank, Adenilton J. da Silva, Lucas P. de Albuquerque, Francesco Petruccione, and Daniel K. Park, ``Compact quantum kernel-based binary classifier,'' \emph{Quantum Science and Technology}, vol.~7, no.~4, Art. no.~045007, 2022, doi: 10.1088/2058-9565/ac842c.

\bibitem{schuld2018book}
Maria Schuld and Francesco Petruccione, \emph{Supervised Learning with Quantum Computers}. Springer, 2018.

\bibitem{plesch2011quantum}
Martin Plesch and \v{C}aslav Brukner, ``Quantum-state preparation with universal gate sets,'' \emph{Physical Review A}, vol.~83, no.~3, Art. no.~032302, 2011, doi: 10.1103/PhysRevA.83.032302.

\bibitem{pcam_01}
Reek Majumdar, Biswaraj Baral, Bhavika Bhalgamiya, and Taposh Dutta Roy, ``Histopathological cancer detection using hybrid quantum computing,'' arXiv:2302.04633v1, 2023.

\bibitem{pcam_03}
N. Goyal, ``Quantum enhanced histopathologic cancer detection,'' M.S. thesis, Dept. Elect. Eng., Arizona State University, Tempe, AZ, USA, Dec. 2024.

\bibitem{pcam_veeling2018}
Bastiaan S. Veeling, Jasper Linmans, Jim Winkens, Taco Cohen, and Max Welling, ``Rotation equivariant CNNs for digital pathology,'' in \emph{Medical Image Computing and Computer Assisted Intervention -- MICCAI 2018}, 21st International Conference, Granada, Spain, Proceedings, Part II, pp.~210--218. Springer International Publishing, 2018.

\bibitem{cover1965geometrical}
Thomas M. Cover, ``Geometrical and statistical properties of systems of linear inequalities with applications in pattern recognition,'' \emph{IEEE Transactions on Electronic Computers}, vol.~EC-14, no.~3, pp.~326--334, 1965, doi: 10.1109/PGEC.1965.264137.

\bibitem{ekert2002direct}
Artur K. Ekert, Carlos M. Alves, Daniel K. L. Oi, Marcin Horodecki, Pawe{\l} Horodecki, and L. C. Kwek, ``Direct estimations of linear and nonlinear functionals of a quantum state,'' \emph{Physical Review Letters}, vol.~88, no.~21, Art. no.~217901, 2002, doi: 10.1103/PhysRevLett.88.217901.

\bibitem{nielsen2000qcqi}
Michael A. Nielsen and Isaac L. Chuang, \emph{Quantum Computation and Quantum Information}. Cambridge University Press, 2000.

\bibitem{preskill2018nisq}
John Preskill, ``Quantum Computing in the NISQ Era and Beyond,'' \emph{Quantum}, vol.~2, p.~79, 2018, doi: 10.22331/q-2018-08-06-79.

\end{thebibliography}

\end{document}
